% Options for packages loaded elsewhere
\PassOptionsToPackage{unicode}{hyperref}
\PassOptionsToPackage{hyphens}{url}
%
\documentclass[
  11pt,
]{article}
\usepackage{amsmath,amssymb}
\usepackage{iftex}
\ifPDFTeX
  \usepackage[T1]{fontenc}
  \usepackage[utf8]{inputenc}
  \usepackage{textcomp} % provide euro and other symbols
\else % if luatex or xetex
  \usepackage{unicode-math} % this also loads fontspec
  \defaultfontfeatures{Scale=MatchLowercase}
  \defaultfontfeatures[\rmfamily]{Ligatures=TeX,Scale=1}
\fi
\usepackage{lmodern}
\ifPDFTeX\else
  % xetex/luatex font selection
\fi
% Use upquote if available, for straight quotes in verbatim environments
\IfFileExists{upquote.sty}{\usepackage{upquote}}{}
\IfFileExists{microtype.sty}{% use microtype if available
  \usepackage[]{microtype}
  \UseMicrotypeSet[protrusion]{basicmath} % disable protrusion for tt fonts
}{}
\makeatletter
\@ifundefined{KOMAClassName}{% if non-KOMA class
  \IfFileExists{parskip.sty}{%
    \usepackage{parskip}
  }{% else
    \setlength{\parindent}{0pt}
    \setlength{\parskip}{6pt plus 2pt minus 1pt}}
}{% if KOMA class
  \KOMAoptions{parskip=half}}
\makeatother
\usepackage{xcolor}
\usepackage[margin=2.2cm]{geometry}
\setlength{\emergencystretch}{3em} % prevent overfull lines
\providecommand{\tightlist}{%
  \setlength{\itemsep}{0pt}\setlength{\parskip}{0pt}}
\setcounter{secnumdepth}{-\maxdimen} % remove section numbering
\ifLuaTeX
  \usepackage{selnolig}  % disable illegal ligatures
\fi
\IfFileExists{bookmark.sty}{\usepackage{bookmark}}{\usepackage{hyperref}}
\IfFileExists{xurl.sty}{\usepackage{xurl}}{} % add URL line breaks if available
\urlstyle{same}
\hypersetup{
  pdftitle={Worksheet 4C: Advanced Rules --- Derivatives and Their Integral Inverses},
  hidelinks,
  pdfcreator={LaTeX via pandoc}}

\title{Worksheet 4C: Advanced Rules --- Derivatives and Their Integral
Inverses}
\author{}
\date{}

\begin{document}
\maketitle

\begin{quote}
\textbf{What This Worksheet Is About}

Each derivative rule has an integral counterpart: - \textbf{Product
Rule} ↔ \textbf{Integration by Parts}\\
- \textbf{Chain Rule} ↔ \textbf{Substitution (u-sub)}

These are inverse operations!
\end{quote}

\begin{center}\rule{0.5\linewidth}{0.5pt}\end{center}

\hypertarget{part-a-the-product-rule-12-problems}{%
\subsection{Part A: The Product Rule (12
problems)}\label{part-a-the-product-rule-12-problems}}

\[\frac{d}{dx}[u \cdot v] = u' \cdot v + u \cdot v'\]

\textbf{Find the derivative:}

\textbf{1.} \(f(x) = x^2(x + 1)\)

\textbf{2.} \(g(t) = t^3(t - 4)\)

\textbf{3.} \(h(p) = (p + 2)(p - 5)\)

\textbf{4.} \(y = (2x + 1)(3x - 2)\)

\textbf{5.} \(f(m) = m^2(m^2 + 3m - 1)\)

\textbf{6.} \(g(r) = (r^2 - 1)(r^2 + 1)\)

\textbf{7.} Verify problem 3 by expanding first.

\textbf{8.} Verify problem 6 by simplifying first:
\((r^2-1)(r^2+1) = ?\)

\begin{center}\rule{0.5\linewidth}{0.5pt}\end{center}

\textbf{With Chain Rule:}

\textbf{9.} \(h(x) = x^2(x + 1)^3\)

\textbf{10.} \(y = (t + 1)^2(t - 1)^3\)

\textbf{11.} \(f(k) = k(2k + 1)^4\)

\textbf{12.} \(g(n) = (n^2 + 1)^2(n - 1)\)

\begin{center}\rule{0.5\linewidth}{0.5pt}\end{center}

\hypertarget{part-b-integration-by-parts-reversing-product-rule-10-problems}{%
\subsection{Part B: Integration by Parts --- Reversing Product Rule (10
problems)}\label{part-b-integration-by-parts-reversing-product-rule-10-problems}}

\[\int u \, dv = uv - \int v \, du\]

\textbf{Connection:} If \(\frac{d}{dx}[f \cdot g] = f'g + fg'\), then
\(\int (f'g + fg') dx = fg + C\)

\begin{center}\rule{0.5\linewidth}{0.5pt}\end{center}

\textbf{13.} The derivative of \(x^2(x+1)\) is \(3x^2 + 2x\).

Therefore: \(\int (3x^2 + 2x) dx ={}\) \_\_\_\_\_\_\_\_\_\_

\textbf{14.} The derivative of \((x-1)(x+3)\) is \(2x + 2\).

Therefore: \(\int (2x + 2) dx ={}\) \_\_\_\_\_\_\_\_\_\_

\textbf{15.} Find \(\frac{d}{dx}[(2x+1)(x-2)]\). Then verify
\(\int [\text{your answer}] dx\) gives back \((2x+1)(x-2) + C\).

\textbf{16.} Start with \(f(x) = x(x+1)^2\). Differentiate, then
integrate your answer.

\begin{center}\rule{0.5\linewidth}{0.5pt}\end{center}

\textbf{Direct Integration (for comparison):}

\textbf{17.} \(\int (3x^2 + 2x) dx ={}\)

\textbf{18.} \(\int (2x + 2) dx ={}\)

\textbf{19.} Compare: Are problems 13 \& 17 the same? Are 14 \& 18 the
same (up to constants)?

\textbf{20.} Explain why integration is ``harder'' than differentiation:

\begin{itemize}
\tightlist
\item
  Differentiation: clear rules always work
\item
  Integration: may require \_\_\_\_\_\_\_\_\_\_ or may not have closed
  form
\end{itemize}

\begin{center}\rule{0.5\linewidth}{0.5pt}\end{center}

\textbf{Preview: Integration by Parts with transcendentals (Worksheet
4D)}

\textbf{21.} We'll need by parts for \(\int x \cdot e^x dx\) because we
can't \_\_\_\_\_\_\_\_\_\_.

\textbf{22.} And for \(\int x \cdot \sin(x) dx\) because
\_\_\_\_\_\_\_\_\_\_.

\begin{center}\rule{0.5\linewidth}{0.5pt}\end{center}

\hypertarget{part-c-the-quotient-rule-10-problems}{%
\subsection{Part C: The Quotient Rule (10
problems)}\label{part-c-the-quotient-rule-10-problems}}

\[\frac{d}{dx}\left[\frac{u}{v}\right] = \frac{u'v - uv'}{v^2}\]

\textbf{Find the derivative:}

\textbf{23.} \(f(x) = \frac{x}{x + 2}\)

\textbf{24.} \(g(t) = \frac{t^2}{t - 1}\)

\textbf{25.} \(h(p) = \frac{p + 1}{p - 1}\)

\textbf{26.} \(y = \frac{x^2 + 1}{x}\)

\textbf{27.} \(f(m) = \frac{1}{m^2 + 1}\)

\textbf{28.} \(g(r) = \frac{r^3}{r^2 + 1}\)

\begin{center}\rule{0.5\linewidth}{0.5pt}\end{center}

\textbf{Quotient Rule vs.~Rewriting:}

\textbf{29.} Find \(\frac{d}{dx}\left[\frac{1}{x^3}\right]\) two ways: -
(a) Quotient rule - (b) Rewrite as \(x^{-3}\)

\textbf{30.} Find \(\frac{d}{dx}\left[\frac{x^4 + 2x^2}{x^2}\right]\) by
simplifying first.

\textbf{31.} Simplify first, then differentiate:
\(\frac{x^2 - 1}{x + 1}\)

\textbf{32.} When MUST you use quotient rule? (When denominator has
\_\_\_\_\_\_\_\_ that don't cancel)

\begin{center}\rule{0.5\linewidth}{0.5pt}\end{center}

\hypertarget{part-d-the-chain-rule-16-problems}{%
\subsection{Part D: The Chain Rule (16
problems)}\label{part-d-the-chain-rule-16-problems}}

\[\frac{d}{dx}[f(g(x))] = f'(g(x)) \cdot g'(x)\]

\textbf{``Derivative of outside × derivative of inside''}

\begin{center}\rule{0.5\linewidth}{0.5pt}\end{center}

\textbf{Find the derivative:}

\textbf{33.} \(f(x) = (x + 3)^4\)

\textbf{34.} \(g(t) = (2t - 1)^5\)

\textbf{35.} \(h(p) = (p^2 + 1)^3\)

\textbf{36.} \(y = (3x^2 - x)^4\)

\textbf{37.} \(f(m) = (m^3 + m)^2\)

\textbf{38.} \(g(r) = (1 - r^2)^6\)

\begin{center}\rule{0.5\linewidth}{0.5pt}\end{center}

\textbf{Negative and Fractional Powers:}

\textbf{39.} \(h(x) = (x + 1)^{-1}\)

\textbf{40.} \(y = (t^2 + 4)^{-2}\)

\textbf{41.} \(f(k) = \sqrt{k + 5}\)

\textbf{42.} \(g(n) = \sqrt{n^2 + 1}\)

\begin{center}\rule{0.5\linewidth}{0.5pt}\end{center}

\textbf{Verify by Expanding:}

\textbf{43.} \((x + 2)^3\): Use chain rule, then expand and
differentiate. Same answer?

\textbf{44.} \((2x - 1)^2\): Both methods.

\begin{center}\rule{0.5\linewidth}{0.5pt}\end{center}

\textbf{Combining Rules:}

\textbf{45.} \(f(x) = x^2(x + 1)^3\) (product + chain)

\textbf{46.} \(g(t) = \frac{t}{(t + 1)^2}\) (quotient + chain)

\textbf{47.} \(h(x) = \sqrt{x(x + 4)}\) (chain with product inside)

\textbf{48.} \(y = \left(\frac{x}{x+1}\right)^3\) (chain on quotient)

\begin{center}\rule{0.5\linewidth}{0.5pt}\end{center}

\hypertarget{part-e-substitution-reversing-the-chain-rule-16-problems}{%
\subsection{Part E: Substitution --- Reversing the Chain Rule (16
problems)}\label{part-e-substitution-reversing-the-chain-rule-16-problems}}

\[\int f(g(x)) \cdot g'(x) \, dx = F(g(x)) + C\]

\textbf{Let \(u = g(x)\), then \(du = g'(x) dx\)}

\begin{center}\rule{0.5\linewidth}{0.5pt}\end{center}

\textbf{Perfect Cases:}

\textbf{49.} \(\int 3x^2(x^3 + 1)^4 \, dx\)

Let \(u = x^3 + 1\), \(du = 3x^2 dx\)

\(= \int u^4 du ={}\)

\textbf{50.} \(\int 2t(t^2 - 3)^5 \, dt\)

\textbf{51.} \(\int 4p^3(p^4 + 2)^{-2} \, dp\)

\textbf{52.} \(\int \frac{6x^2}{\sqrt{x^3 + 1}} \, dx\)

\textbf{53.} \(\int (2m + 1)(m^2 + m)^3 \, dm\)

\begin{center}\rule{0.5\linewidth}{0.5pt}\end{center}

\textbf{Adjusting Constants:}

\textbf{54.} \(\int x(x^2 + 1)^3 \, dx\)

We need \(2x dx\) but have \(x dx\). So:
\(x dx = \frac{1}{2}(2x dx) = \frac{1}{2}du\)

\(= \frac{1}{2}\int u^3 du ={}\)

\textbf{55.} \(\int t(t^2 + 5)^4 \, dt\)

\textbf{56.} \(\int \frac{x}{\sqrt{x^2 + 4}} \, dx\)

\textbf{57.} \(\int p^2(p^3 - 1)^2 \, dp\)

\begin{center}\rule{0.5\linewidth}{0.5pt}\end{center}

\textbf{Verify by Differentiating:}

\textbf{58.} Check problem 49: Differentiate \(\frac{(x^3+1)^5}{5}\)
using chain rule.

\textbf{59.} Check problem 54: Differentiate \(\frac{(x^2+1)^4}{8}\).

\begin{center}\rule{0.5\linewidth}{0.5pt}\end{center}

\textbf{Connection to Chain Rule:}

\textbf{60.} Chain rule: \(\frac{d}{dx}[(x^2+1)^5] = 10x(x^2+1)^4\)

Therefore: \(\int 10x(x^2+1)^4 dx ={}\) \_\_\_\_\_\_\_\_\_\_

\textbf{61.} Chain rule: \(\frac{d}{dx}[(t^3-2)^4] = 12t^2(t^3-2)^3\)

Therefore: \(\int 12t^2(t^3-2)^3 dt ={}\) \_\_\_\_\_\_\_\_\_\_

\textbf{62.} Explain: Substitution is ``chain rule backwards.''

\begin{center}\rule{0.5\linewidth}{0.5pt}\end{center}

\textbf{More Practice:}

\textbf{63.} \(\int (3x - 2)^4 \, dx\)

\textbf{64.} \(\int (5t + 1)^{-3} \, dt\)

\begin{center}\rule{0.5\linewidth}{0.5pt}\end{center}

\hypertarget{part-f-combining-everything-10-problems}{%
\subsection{Part F: Combining Everything (10
problems)}\label{part-f-combining-everything-10-problems}}

\textbf{Differentiate:}

\textbf{65.} \(\frac{x^2}{(x+1)^3}\)

\textbf{66.} \(x(x^2 + 1)^{1/2}\)

\textbf{67.} \(\frac{(x-1)^2}{(x+1)^2}\)

\begin{center}\rule{0.5\linewidth}{0.5pt}\end{center}

\textbf{Integrate:}

\textbf{68.} \(\int \frac{2x}{(x^2 + 1)^3} dx\)

\textbf{69.} \(\int \frac{x}{\sqrt{x^2 + 9}} dx\)

\textbf{70.} \(\int (x+1)(x^2 + 2x)^4 dx\)

\emph{Hint: \(\frac{d}{dx}(x^2 + 2x) = 2x + 2 = 2(x + 1)\)}

\begin{center}\rule{0.5\linewidth}{0.5pt}\end{center}

\textbf{Round-Trips:}

\textbf{71.} Start with \((x^2 + 1)^3\). Differentiate, then integrate
back.

\textbf{72.} Start with \(x(x + 1)^2\). Differentiate, then integrate
back.

\textbf{73.} Start with \(\frac{x}{x+1}\). Differentiate, then integrate
back.

\textbf{74.} Why does ``differentiate then integrate'' give
\(f(x) + C\), but ``integrate then differentiate'' gives exactly
\(f(x)\)?

\begin{center}\rule{0.5\linewidth}{0.5pt}\end{center}

\end{document}
